\section{Capacity-controlled signed data in the strong energy dual}
\label{sec:r12-capacity-dual-data}

We retain exactly the strict Laurent-box and extension hypotheses
(ED1), the arbitrary original middle-positive separately closed
background $S$, and all original tests. Put $d=q+1\geq2$,
$\mathscr C=\widehat{\mathcal E}^{1}_{\rm str}(V)$, and
$\mathbb D_S=\mathscr C-\mathscr C$. No moment presentation or
special form of $S$ is imposed in this section. The outer capacity
$C_S^*(E;V)$ always uses the full original bounded classical tester
class on the specified whole domain $V$.

The datum criterion below is a sufficient condition for membership
in the strong-continuous dual, not a nonlinear solvability theorem.
It does not require domination by one actual curvature measure.

\begin{definition}[An integrable capacity modulus]
\label{def:r12-integrable-capacity-modulus}
An integrable capacity modulus in degree $d$ is a nondecreasing
function $F:[0,\infty]\to[0,\infty]$, finite on every finite
positive argument, with $F(0)=0$, such that
\[
                 \int_0^1 F(s)s^{-1-1/d}\,ds<\infty.
 \tag{CD1}
\]
The value at infinity is its increasing limit. Let $\mu$ be a
finite signed original Radon measure on $V$. We say that its
variation has modulus $F$ if
\[
                    |\mu|(E)\leq F\bigl(C_S^*(E;V)\bigr)
                    \quad\text{for every Borel }E\subset V.
 \tag{CD2}
\]
The total variation $m=|\mu|(V)$ is finite independently of
whether $C_S^*(V;V)$ is finite.
\end{definition}

Put $C_0=d2^d$, as in (ED5), and define
\[
       A_{F,m}:=\int_0^\infty
              \min\{m,F(C_0s^{-d})\}\,ds.
 \tag{CD3}
\]
If $m=0$, set $A_{F,0}=0$. Substitution $x=C_0s^{-d}$ on a
tail of the integral shows that (CD1) makes (CD3) finite.

\begin{theorem}[Literal signed data and strong continuity]
\label{thm:r12-capacity-modulus-strong-dual}
Suppose $\mu$ satisfies (CD2) for an integrable capacity modulus.
Then every $h\in\mathbb D_S$ is absolutely integrable for
$|\mu|$, and
\[
       \ell_\mu(h)=-\int_V h\,d\mu,
       \qquad
       \|\ell_\mu\|_*\leq A_{F,m}
 \tag{CD4}
\]
defines an element of $\mathbb D_{S,\mathrm s}^*$.
More explicitly, if $u_j,u\in\mathscr C$ and $I(u_j,u)\to0$,
then
\[
                    \int_V|u_j-u|\,d|\mu|\longrightarrow0.
 \tag{CD5}
\]
All differences and integrals are unchanged by locally outer-null
modifications. The assertion holds for every $q\geq1$ and every
allowed original $S$, including degenerate backgrounds after the
specified relative quotient.
\end{theorem}

\noindent\textbf{Status: \texttt{NEW\_PROVED}} as a
full-degree signed datum class in the newly constructed Banach
dual. The sublevel and all-local strong capacity estimates are
the explicitly inherited R11 inputs (CI3), (EV7)--(EV8).

\begin{proof}
For $a\in\mathscr C$, (ED5) and (CD2) give
\[
       |\mu|\{a<-t\}
       \leq \min\{m,F(C_0r(a)^d t^{-d})\},\qquad t>0.
 \tag{CD6}
\]
The actual Borel representative of $a$ is nonpositive and finite
outside an outer-null set. Condition (CD2) annihilates this set.
The layer-cake formula, followed by $t=r(a)s$, therefore gives
\[
                        \int_V(-a)\,d|\mu|
                                      \leq A_{F,m}r(a).
 \tag{CD7}
\]
When $r(a)=0$, the class is zero outside a local outer-null set,
which is globally outer-null in the fixed domain by the R11
countable compact exhaustion and domain comparison. Thus (CD7)
holds in that case too. If $h=a-b$, then
$|h|\leq(-a)+(-b)$ off the joint null set. Consequently
$\int|h|\,d|\mu|\leq A_{F,m}(r(a)+r(b))$. Taking the infimum
over its presentations proves absolute integrability and (CD4).
This is a literal integral on the actual Borel difference space;
no unproved extension of a measure to abstract completion points
is used.

Now let $u_j\to u$ strongly. The energy-root estimate gives a
finite $R$ with $r(u_j),r(u)\leq R$. If $R=0$ the assertion
is immediate. Otherwise, for every $L>0$, (CD6) gives
\[
 \sup_j\int_V\bigl((-u_j)-L\bigr)_+\,d|\mu|
       \leq\int_L^\infty
                  \min\{m,F(C_0R^d t^{-d})\}\,dt
                   \longrightarrow0\quad(L\to\infty).
 \tag{CD8}
\]
The same bound holds for $u$. This is a uniform integrability
statement for the full finite datum, including its boundary tail.

Fix $L$ and write $u_j^{(L)}=\max(u_j,-L)$ and
$u^{(L)}=\max(u,-L)$. On each compact $K\Subset V$, the
R11 strong estimate makes
$C_S^*(K\cap\{|u_j^{(L)}-u^{(L)}|>\varepsilon\};V)$ tend
to zero. An integrable nondecreasing modulus satisfies
$F(s)\to0$ as $s\downarrow0$: otherwise (CD1) would diverge.
Hence (CD2) gives convergence in $|\mu|$-measure on $K$.
The clipped differences are bounded by $L$, so their integrals
over $K$ tend to zero. Radon inner regularity and finiteness of
$|\mu|$ allow a compact $K$ with arbitrarily small
$|\mu|(V\setminus K)$; the same bound $L$ then passes this
convergence to the whole domain. Finally (CD8) removes the
clipping uniformly. This proves (CD5) and the defining
strong-dual continuity (NI4).

Every capacity in the argument is the frozen full original
capacity. Actual moment refinements alter neither the measure
integrals nor the classes of tests defining these estimates.
No assertion of raw outer regularity on arbitrary Borel sets
is needed.
\end{proof}

\begin{corollary}[Power and logarithmic endpoint moduli]
\label{cor:r12-power-capacity-dual-data}
Let $m=|\mu|(V)<\infty$. If, for constants $A>0$ and
$\gamma>1/d$,
\[
                 |\mu|(E)\leq A\,C_S^*(E;V)^\gamma
                    \quad(E\text{ Borel}),
 \tag{CD9}
\]
then $\ell_\mu\in\mathbb D_{S,\mathrm s}^*$ and
\[
 \|\ell_\mu\|_*
 \leq \frac{d\gamma}{d\gamma-1}
       m^{1-1/(d\gamma)}
                  (A C_0^\gamma)^{1/(d\gamma)}.
 \tag{CD10}
\]
There are also endpoint admissible moduli: any nondecreasing
extension, for $s$ bounded away from zero, of
\[
             F(s)=A s^{1/d}\bigl(\log(e/s)\bigr)^{-\alpha}
                   \quad(0<s\leq1),\qquad \alpha>1,
 \tag{CD11}
\]
satisfies (CD1). Neither the unmodified critical power
$s^{1/d}$ nor the range $\alpha\leq1$ satisfies this
integral criterion. This last statement is about the criterion,
not a universal necessity claim for all possible measures.
\end{corollary}

\noindent\textbf{Status: \texttt{DIRECT\_COROLLARY}}.

\begin{proof}
For (CD9), (CD3) becomes the integral of
$\min(m,A C_0^\gamma s^{-d\gamma})$. Splitting at
$s_0=(A C_0^\gamma/m)^{1/(d\gamma)}$ yields (CD10) when
$m>0$, and the zero case is immediate. For (CD11), substitution
$y=\log(e/s)$ turns (CD1), up to the harmless constant $A$,
into $\int_1^\infty y^{-\alpha}\,dy$. This proves all printed
assertions without an endpoint counterexample being inferred
from the failure of a sufficient estimate.
\end{proof}

\begin{proposition}[Uniform capacity control converts variation to dual convergence]
\label{prop:r12-capacity-dini-data-stability}
Let $\mu_j,\mu$ be finite signed Radon measures such that
$m_j=|\mu_j-\mu|(V)\to0$. Suppose the variations
$|\mu_j-\mu|$ all satisfy (CD2) with one common integrable
modulus $F$. Then
\[
         \|\ell_{\mu_j-\mu}\|_*
                \leq A_{F,m_j}\longrightarrow0.
 \tag{CD12}
\]
This unconditional statement defines only the functional of the
controlled difference. If the individual datum functionals are
defined, it is their difference by literal integration.
If additionally each $\ell_{\mu_j}$ and $\ell_\mu$ lies in
$\mathbb D_{S,\mathrm s}^*$, their canonical variational
potentials from (NI7) converge strongly, and their normal
residuals converge in dual norm. For power domination the
right side in (CD12) has the explicit bound (CD10) with
$m=m_j$. If $\ell_{\mu_j}=J(u_j)$ are actual completed source
functionals, then the limit is itself $J(u)$; in that special
case $\widehat\mu_u=\mu$ on every original scalar test.
\end{proposition}

\noindent\textbf{Status: \texttt{DIRECT\_COROLLARY}} of
the new datum criterion and nonlinear closed-image theorem.

\begin{proof}
Apply (CD4) to $\mu_j-\mu$. Since $m_j$ is bounded, the
integrands in (CD3) are dominated by one integrable function
$\min(M,F(C_0s^{-d}))$ and tend pointwise to zero. Dominated
convergence proves (CD12). The quantitative retraction estimate
(NI11) supplies strong convergence of the variational potentials
and dual convergence of their curvatures; subtracting from the
data gives convergence of the residuals. The final assertion is
the closed-image criterion (NI1)--(NI3). It assumes source
approximants and does not create them from capacity domination.
\end{proof}

\paragraph{Relation to the frozen first-order theory.}
R11 already solves a different linear transposition equation for
arbitrary finite diffuse data in a first-order Dirichlet-form
space. That result does not say that the solution is a classical
psh strong-completion member. The present theorem places signed
data in an all-degree nonlinear energy dual and supplies its
variational decomposition. It neither supersedes that linear
result nor removes the nonlinear normal residual. In particular,
bounded-source domination implies (CD9) with $\gamma=1$, but
the compact-source examples in Section~\ref{sec:r12-positive-moment-lifts}
show that even this datum class does not imply the unrestricted
source equation.
